\subsection{基线方法}
\label{sec:baselines}

为对所提 PMDF-Net 进行上下文化比较，我们对比了涵盖经典信号处理和现代深度学习的五种成熟去噪方法。所有基线方法在比较前经过仔细调优以获得最佳性能。

\textbf{信号平均}代表激光超声检测中信噪比(SNR)改善的传统工业标准。获取重复 N 次平均信号并与 M 次参考（其中 M$\gg$N）进行比较。预期信噪比(SNR)改善规模为$10\log_{10}(M/N)$dB，尽管这以延长激光曝光和降低检测吞吐量为代价。该基线直接捕捉了本研究旨在克服的权衡。

\textbf{维纳滤波}采用自适应方法，从输入信号估计噪声功率谱密度并相应地应用频域滤波\cite{dabov2007bm3d}。与简单平均不同，维纳滤波利用信号和噪声分量之间的频谱差异。该方法需要准确的噪声表征；当该假设成立时，它在不增加多次采集时间成本的情况下提供中等去噪效果。

\textbf{BM3D}（块匹配 3D 协同滤波）通过将相似信号块分组并在稀疏 3D 表示中应用协同滤波来扩展常规 3D 变换域去噪\cite{dabov2007bm3d}。BM3D 在自然图像上取得优异性能，并已适应 1D 信号去噪。其优势在于利用块相似性；
然而，它可能在激光超声波形的瞬态、高频特征上表现困难，因为精确块匹配较少见。

\textbf{小波去噪（DWT）}使用 db4 小波族通过软阈值化应用离散小波变换\cite{chen2019wavelet}。小波阈值水平通过在留置噪声数据上的交叉验证来选择，以平衡噪声抑制与信号失真。
这种方法提供了一个计算效率高的基线，捕获多尺度信号特征，尽管固定小波基选择可能无法最佳地表示实践中遇到的多样声波形。

\textbf{单域 CNN（DnCNN/U-Net）}仅在时域中训练深度卷积去噪网络\cite{zhang2017denoising}。该消融基线反映了 PMDF-Net 的时域分支，但没有时频分支或物理约束。
它用于隔离多域融合的贡献，预期显示与 PMDF-Net 可比的信噪比(SNR)改善，但由于缺乏频谱引导和基于物理的损失项而特征保留下降。

每个基线通过在相关超参数上（如维纳噪声方差估计窗口、BM3D 匹配阈值、DWT 分解层、CNN 学习率和轮次）进行网格搜索调优以确保公平比较。